% mathtools documentation
\documentclass{article}

\usepackage{fourier}
\usepackage{mathtools}

\usepackage{xcolor,varioref,amssymb}

\usepackage{xurl}

\usepackage[final,
hyperindex=false,
colorlinks,
]{hyperref}

\makeatletter
\newcommand*\thinfbox[2][black]{\fboxsep0pt\textcolor{#1}{\rulebox{{\normalcolor#2}}}}
\newcommand*\thinboxed[2][black]{\thinfbox[#1]{\ensuremath{\displaystyle#2}}}
\newcommand*\rulebox[1]{%
  \sbox\z@{\ensuremath{\displaystyle#1}}%
  \@tempdima\dp\z@
  \hbox{%
    \lower\@tempdima\hbox{%
      \vbox{\hrule height\fboxrule\box\z@\hrule height\fboxrule}%
    }%
  }%
}
\makeatother

\title{BookML test: \texttt{mathtools} user guide}
\author{Morten Høgholm, Lars Madsen and the \LaTeX3 project}
\date{2026-03-02}

\usepackage{bookml/bookml}
\usepackage{framed}
\AtBeginDocument{\begin{framed}
  \begin{description}
    \item[BookML] \BookMLversion
    \item[\LaTeXML] \LaTeXMLversion{} (\csname bml@generator@engine\endcsname)
    \item[\LaTeX] \fmtversion
  \end{description}
\end{framed}}


\begin{document}

\section{Tools for mathematical typesetting}

\subsubsection{A complement to \texttt{\textbackslash smash},
  \texttt{\textbackslash llap}, and \texttt{\textbackslash rlap}}

\[
  X = \sum_{1\le i\le j\le n} X_{ij}
\]
\[
  X = \sum_{\mathclap{1\le i\le j\le n}} X_{ij}
\]

\subsubsection{Forcing a cramped style}

\begin{equation}\label{eq:mathclap}
  \sum_{\mathclap{a^2<b^2<c}}\qquad
  \sum_{a^2<b^2<c}
\end{equation}

\begingroup \fontsize{24}{\baselineskip}\selectfont
\[
  \sum_{\mathclap{a^2<b^2<c}}\qquad
  \sum_{a^2<b^2<c}
\]
\endgroup

\[
  \cramped{x^2}               \leftrightarrow x^2    \quad
  \cramped[\scriptstyle]{x^2} \leftrightarrow {\scriptstyle x^2}
\]

\begin{quote}
  The 2005 Euro\TeX{} conference is held in Abbaye des
  Pr\'emontr\'es, France, marking the 16th ($2^{2^2}$) anniversary
  of both Dante and GUTenberg (the German and French \TeX{} users
  group resp.).
\end{quote}


\begin{quote}
  The 2005 Euro\TeX{} conference is held in Abbaye des
  Pr\'emontr\'es, France, marking the 16th ($\cramped{2^{2^2}}$)
  anniversary of both Dante and GUTenberg (the German and French
  \TeX{} users group resp.).
\end{quote}

\begin{equation*}\label{eq:mathclap-b}
  \sum_{\crampedclap{a^2<b^2<c}}
  \tag{\ref{eq:mathclap}*}
\end{equation*}

\begingroup \fontsize{24}{\baselineskip}\selectfont
\[
  \sum_{\substack{a^2<b^2<c}}\qquad
  \sum_{a^2<b^2<c}
\]
\endgroup

\begingroup \fontsize{24}{\baselineskip}\selectfont
\[
  \sum_{\crampedsubstack{a^2<b^2<c}}\qquad
  \sum_{a^2<b^2<c}
\]
\endgroup

\subsubsection{Smashing an operator}
\[
  V = \sum_{1\le i\le j\le n}^{\infty} V_{ij}                  \quad
  X = \smashoperator{\sum_{1\le i\le j\le n}^{3456}} X_{ij}    \quad
  Y = \smashoperator[r]{\sum\limits_{1\le i\le j\le n}} Y_{ij} \quad
  Z = \smashoperator[l]{\mathop{T}_{1\le i\le j\le n}} Z_{ij}
\]

\subsubsection{Adjusting limits of operators}

\[
  \text{a)} \lim_{n\to\infty} \max_{p\ge n} \quad
  \text{b)} \lim_{n\to\infty} \max_{p^2\ge n} \quad
  \text{c)} \lim_{n\to\infty} \sup_{p^2\ge nK} \quad
  \text{d)} \limsup_{n\to\infty} \max_{p\ge n}
\]

\[
  \text{a)} \adjustlimits\lim_{n\to\infty} \max_{p\ge n} \quad
  \text{b)} \adjustlimits\lim_{n\to\infty} \max_{p^2\ge n} \quad
  \text{c)} \adjustlimits\lim_{n\to\infty} \sup_{p^2\ge nK} \quad
  \text{d)} \adjustlimits\limsup_{n\to\infty} \max_{p\ge n}
\]


\subsubsection{Swapping space above \texorpdfstring{\AmS}{AMS} display math environments }

\noindent\rule\textwidth{1pt}
\begin{align*}  A &= B \end{align*}
\noindent\rule\textwidth{1pt}
\begin{align*}
\SwapAboveDisplaySkip
A &= B
\end{align*}

\subsection{Controlling tags}

\newtagform{brackets}{[}{]}
\usetagform{brackets}
  \begin{equation}
    E \neq m c^3
  \end{equation}

\newtagform{brackets2}[\textbf]{[}{]}
\usetagform{brackets2}
\begin{equation}
  E \neq m c^3
\end{equation}

\subsubsection{Showing only referenced tags}

\begin{quote}\renewcommand*\rmdefault{ppl}\normalfont\itshape
\begin{equation*}
  a=b \label{eq:example}\tag*{Q\&A}
\end{equation*}
See \ref{eq:example} or is it better with \refeq{eq:example}?
\end{quote}

\mathtoolsset{showonlyrefs,showmanualtags}
\usetagform{brackets}
\begin{gather}
  a=a \label{eq:a} \\
  b=b \label{eq:b} \tag{**}
\end{gather}
This should refer to the equation containing $a=a$: \eqref{eq:a}.
Then a switch of tag forms.
\usetagform{default}
\begin{align}
  c&=c \label{eq:c} \\
  d&=d \label{eq:d}
\end{align}
This should refer to the equation containing $d=d$: \eqref{eq:d}.
\begin{equation}
  e=e
\end{equation}
Back to normal.\mathtoolsset{showonlyrefs=false}
\begin{equation}
  f=f
\end{equation}

\subsection{Extensible symbols}

\[
  A \xLeftarrow[under]{over} B
\]

\begin{align*}
  A &\xleftrightarrow[below]{above} B &
  A &\xRightarrow[below]{above} B \\
  A &\xLeftarrow[below]{above} B &
  A &\xLeftrightarrow[below]{above} B \\
  A &\xhookleftarrow[below]{above} B &
  A &\xhookrightarrow[below]{above} B \\
  A &\xmapsto[below]{above} B \\
  A &\xrightharpoondown[below]{above} B &
  A &\xrightharpoonup[below]{above} B \\
  A &\xleftharpoondown[below]{above} B &
  A &\xleftharpoonup[below]{above} B \\
  A &\xrightleftharpoons[below]{above} B &
  A &\xleftrightharpoons[below]{above} B \\
\end{align*}

\begin{align*}
  A & \longrightarrow B & A & \longleftarrow B \\
  A & \xlongrightarrow[b]{a} B & A & \xlongleftarrow[b]{a} B \\
  A & \xlongrightarrow[below]{above} B & A & \xlongleftarrow[below]{above} B \\
\end{align*}
\begin{align*}
  A & \xLongrightarrow B & A & \Longleftarrow B \\
  A & \xLongrightarrow[b]{a} B & A & \xLongleftarrow[b]{a} B \\
  A & \xLongrightarrow[below]{above} B & A & \xLongleftarrow[below]{above} B \\
\end{align*}

\subsubsection{Braces and brackets}

\[
  \underbracket[3pt]{xxx\  yyy}_{zzz} \quad \text{and} \quad
  \underbracket[1pt][7pt]{xxx\  yyy}_{zzz}
\]
\[
  \mathop{\thinboxed[blue]{\sum}}_{n}
  \mathop{\thinboxed[blue]{\LaTeXunderbrace{\thinboxed[green]{foof}}}}_{zzz}
\]

\par\Huge\vskip-\baselineskip
\[
  \mathop{\thinboxed[blue]{\sum}}_{n}
  \mathop{\thinboxed[blue]{\LaTeXunderbrace{\thinboxed[green]{foof}}}}_{zzz}
\]\normalsize
\par\footnotesize
\[
  \mathop{\thinboxed[blue]{\sum}}_{n}
  \mathop{\thinboxed[blue]{\LaTeXunderbrace{\thinboxed[green]{foof}}}}_{zzz}
\]\normalsize

\subsection{New mathematical building blocks}

\subsubsection{Matrices}\label{subsubsec:matrices}

\[
  \begin{pmatrix*}[r]
    -1 & 3 \\
    2  & -4
  \end{pmatrix*}
\]

\[
\begin{bsmallmatrix}    a & -b \\ -c & d \end{bsmallmatrix}
\begin{bsmallmatrix*}[r] a & -b \\ -c & d \end{bsmallmatrix*}
\]
\[
\begin{bsmallmatrix} a & -b \\ -c & d \end{bsmallmatrix}
\begin{bsmallmatrix*}[r] a & -b \\ -c & d \end{bsmallmatrix*}
\]

\subsubsection{The \texttt{multlined} environment}

\[
  A = \begin{multlined}[t]
        \framebox[4cm]{first} \\
        \framebox[4cm]{last}
      \end{multlined} B
\]
\[
  \begin{multlined}
    \framebox[.65\columnwidth]{First line} \\
    \framebox[.5\columnwidth]{Second line} \\
    \shoveleft{L+E+F+T}         \\
    \shoveright{R+I+G+H+T}         \\
    \shoveleft[1cm]{L+E+F+T}         \\
    \shoveright[\widthof{$R+I+G+H+T$}]{R+I+G+H+T}         \\
    \framebox[.65\columnwidth]{Last line}
  \end{multlined}
\]

\[
    \begin{multlined}[b][7cm]
        \framebox[4cm]{first} \\
        \framebox[4cm]{last}
      \end{multlined} = B
\]

\subsubsection{More \texttt{cases}-like environments}

\[
  a=\begin{cases}
    E = m c^2     & \text{Nothing to see here} \\
    \int x-3\, dx & \text{Integral is text style}
  \end{cases}
\]

\[
  \begin{dcases}
    E = m c^2     & c \approx 3.00\times 10^{8}\,\mathrm{m}/\mathrm{s} \\
    \int x-3\, dx & \text{Integral is display style}
  \end{dcases}
\]

\[
  a= \begin{dcases*}
    E = m c^2     & Nothing to see here \\
    \int x-3\, dx & Integral is display style
  \end{dcases*}
\]

\[
\begin{rcases*}
  x^2 & for $x>0$\\
  x^3 & else
\end{rcases*} \quad \Rightarrow\cdots
\]


\subsubsection{Emulating indented lines in alignments}
\begin{align*}
  \MoveEqLeft \framebox[10cm][c]{Long first line}\\
  & = \framebox[6cm][c]{2nd line}\\
  & \leq \dots
\end{align*}

\begin{align*}
  \MoveEqLeft[3] \framebox[10cm][c]{Long first line}\\
    = {} & \framebox[6cm][c]{2nd line}\\
        & + \framebox[4cm][c]{last part}
\end{align*}

\subsubsection{Boxing a single line in an alignment}

\begin{align*}
  \Aboxed{ f(x) & = \int h(x)\, dx} \\
                & = g(x)
\end{align*}

\begin{align*}
  \setlength\fboxsep{1em}
  \Aboxed{ f(x) &= 0 } & \Aboxed{ g(x) &= b} \\
  \Aboxed{ h(x) }      & \Aboxed{ i(x) }
\end{align*}

\newcommand\myboxA[1]{\fcolorbox{gray}{yellow}{#1}}
\newcommand\myboxB[1]{{\fboxsep=5pt\fboxrule=2pt\fbox{#1}}}
\MakeAboxedCommand\AboxedA\myboxA
\MakeAboxedCommand\AboxedB\myboxB
\begin{align*}
  \AboxedA{ f(x) &= x^2 } & \AboxedB{ f(x) &= x^2 }
\end{align*}

\subsubsection{Adding arrows between lines in an alignment}

\begin{alignat}{2}
  && \framebox[1.5cm]{} &= \framebox[3cm]{}\\
  \ArrowBetweenLines
  && \framebox[1.5cm]{} &= \framebox[2cm]{}
\end{alignat}

\begin{alignat*}{2}
  \framebox[1.5cm]{} &= \framebox[3cm]{}  &&\\
  \ArrowBetweenLines*[\Downarrow]
  \framebox[1.5cm]{} &= \framebox[2cm]{}  &&
\end{alignat*}

\subsubsection{Centered \texorpdfstring{\texttt{vdots}}{\textbackslash vdots}}

\begin{align*}
  \framebox[1.5cm]{} &= \framebox[3cm]{}\\
                      & \vdots\\
                      &= \framebox[3cm]{}
\end{align*}

\begin{align*}
  a &= b              \\
    & \vdotswithin{=} \\
    & = c             \\
    \shortvdotswithin{=}
    & = d
\end{align*}

\begin{alignat*}{3}
  A&+ B &&= C &&+ D \\
  \MTFlushSpaceAbove
  &\vdotswithin{+} &&&& \vdotswithin{+}
  \MTFlushSpaceBelow
  C &+ D &&= Y &&+K
\end{alignat*}

\subsection{Intertext and short intertext}

\begin{align}
  a&=b \intertext{Some text}
  c&=d
\end{align}

\begin{align}
  a&=b \shortintertext{Some text}
  c&=d
\end{align}

\subsection{Paired delimiters}


\DeclarePairedDelimiter\abs\lvert\rvert
\[
  \abs{\frac{a}{b}}
\]

\[
  \abs*{\frac{a}{b}}
\]

\[
  \abs[\Bigg]{\frac{a}{b}}
\]

\DeclarePairedDelimiterX\innerp[2]{\langle}{\rangle}{#1,#2}
\DeclarePairedDelimiterX\braket[3]{\langle}{\rangle}%
{#1\,\delimsize\vert\,\mathopen{}#2\,\delimsize\vert\,\mathopen{}#3}
\[
\innerp*{A}{ \frac{1}{2} } \quad
\braket[\Big]{B}{\sum_{k}}{C}
\]

\providecommand\given{}
\newcommand\SetSymbol[1][]{\nonscript\:#1\vert\allowbreak\nonscript\:\mathopen{}}
\DeclarePairedDelimiterX\Set[1]\{\}{%
    \renewcommand\given{\SetSymbol[\delimsize]}
    #1
}
\[ \Set*{ x \in X \given \frac{\sqrt{x}}{x^2+1} > 1 } \]

\subsection{Special symbols}

\subsubsection{Vertically centered colon}
\[\mathtoolsset{centercolon}
  a := b
\]

\subsubsection{A few additional symbols}

\[
    \lim_{a\ndownarrow 0} f(a) \neq \bigtimes_n X_n \qquad
    \frac{ \bigtimes_{k=1}^7 B_k \nuparrow \Omega }{2}
\]


\subsection{Mathematics within italic text}

\begin{quote}\itshape
Compare these lines: \par
\mathtoolsset{mathic}
Subset of \(V\) and subset of \(A\). \par
\mathtoolsset{mathic=false}
Subset of \(V\) and subset of \(A\).
\par
\end{quote}

\subsection{Left sub/superscripts}

\[
  {}^{4}_{12}\mathbf{C}^{5+}_{2}          \quad
  \prescript{14}{2}{\mathbf{C}}^{5+}_{2}  \quad
  \prescript{4}{12}{\mathbf{C}}^{5+}_{2}  \quad
  \prescript{14}{}{\mathbf{C}}^{5+}_{2}   \quad
  \prescript{}{2}{\mathbf{C}}^{5+}_{2}
\]

\newcommand*\myisotope[3]{%
  \begingroup
    \mathtoolsset{
      prescript-sup-format=\mathit,
      prescript-sub-format=\mathbf,
      prescript-arg-format=\mathrm,
    }%
  \prescript{#1}{#2}{#3}%
  \endgroup
}
\[
  \myisotope{A}{Z}{X}\to \myisotope{A-4}{Z-2}{Y}+
  \myisotope{4}{2}{\alpha}
\]

\subsection{Spreading equations}\label{sec:spread}

\begin{spreadlines}{20pt}
Large spaces between the lines.
\begin{gather}
  a=b\\
  c=d
\end{gather}
\end{spreadlines}
Back to normal spacing.
\begin{gather}
  a=b\\
  c=d
\end{gather}


\subsection{Gathered environments}\label{subsec:gathered}

\begin{equation}
  \begin{lgathered}
      x=1,\quad x+1=2 \\
      y=2
  \end{lgathered}
\end{equation}

\newcounter{steplinecnt}
\newcommand\stepline{\stepcounter{steplinecnt}\thesteplinecnt}
\newgathered{stargathered}{\llap{\stepline}$*$\quad\hfil}{\hfil}{\setcounter{steplinecnt}{0}}
\begin{gather}
  \begin{stargathered}
    x=1,\quad x+1=2 \\
    y=2
  \end{stargathered}
\end{gather}

\subsection{Split fractions}

\[
a=\frac{
    \splitfrac{xy + xy + xy + xy + xy}
              {+ xy + xy + xy + xy}
  }
  {z}
=\frac{
    \splitdfrac{xy + xy + xy + xy + xy}
              {+ xy + xy + xy + xy}
  }
  {z}
\]
\[
\frac{
    \splitfrac{xy + xy + xy + xy + xy}
    {
    \splitfrac{xy + xy + xy + xy + xy}
              {+ xy + xy + xy + xy}
    }
  }
  {z}
=\frac{
    \splitfrac{xy + xy + xy + xy + xy}
    {
    \splitfrac{\mathstrut  xy + xy + xy + xy + xy}
              {+ xy + xy + xy + xy}
    }
  }
  {z}
\]

\section{New additions}

 \subsection{Variable math strut}
\[
  \begin{cases*}
    \frac{\frac{ x-1 }{ x-\sin x} }{ \sqrt{ 1 -x }} & $x >0$ \\
    0 & otherwise
  \end{cases*}
\qquad\text{vs.}\qquad
  \begin{cases*}
    \frac{\frac{ \xmathstrut{0.1} x-1 }{ \xmathstrut{0.25} x-\sin x} }{\xmathstrut{0.4} \sqrt{ 1 -x }} & $x >0$ \\
    0 & otherwise
  \end{cases*}
\]

\newcommand\vfb[1]{\begingroup\fboxsep=0pt\boxed{\,#1\,}\endgroup}
\[
  \llap{\rlap{\rule{18mm}{0.1pt}}\quad}a \vfb{\mathstrut} \ \vfb{\xmathstrut{0}}\
  \vfb{\xmathstrut{0.5} } \  \vfb{\xmathstrut{-0.1} }\  \vfb{\xmathstrut[0.5]{0}} a
\]

\end{document}
